\section{Coincidence Classes}
\label{sec:Coincidence Classes}
\subsection{Coincidence}
\label{sub:Coincidence}
\begin{frame}{Coincidence}
  \begin{definition}
    Two mesh patterns are said to be \emph{coincident} if they are avoided by the
    same set of permutations at every length.
  \end{definition}
  \begin{block}{}
    Distinct classical patterns can never be coincident.

    Coincidences of mesh patterns of length \(2\) are classified.

    Aim to establish rules that classify coincidences when we
    have one mesh pattern and one classical pattern.

    We call the classical pattern a \emph{dominating pattern}.
  \end{block}
\end{frame}

\subsection{Summary of Experimental Results}
\label{sub:Summary of Experimental Results}
\begin{frame}{Experimental Results}
  \begin{block}{}
  \begin{table}[htb]
\begin{center}
\begin{tabularx}{\textwidth}{c|Y|Y|Y|Y|}
\cline{2-5}
& \multicolumn{4}{|c|}{Dominating Pattern}\\
\cline{2-5}
& \multicolumn{2}{|c|}{\(\perm{2,3,1}\)}&\multicolumn{2}{|c|}{\(\perm{3,2,1}\)}\\
\cline{2-5}
& \(\perm{1,2}\)&\(\perm{2,1}\)&\(\perm{1,2}\)&\(\perm{2,1}\)\\
\hline
\multicolumn{1}{|r|}{No Dominating rule}&\(220\)&\(220\)&\(220\)&\(220\)\\
\hline
\multicolumn{1}{|r|}{\hspace{0pt}\invisible<2->{First Dominating rule}}&\hspace{0pt}\invisible<2->{\(85\)}&\hspace{0pt}\invisible<2->{\(43\)}&\hspace{0pt}\invisible<2->{\(220\)}&\hspace{0pt}\invisible<2->{\(29\)}\\
\hline
\multicolumn{1}{|r|}{\hspace{0pt}\invisible<2->{Second Dominating rule}}&\hspace{0pt}\invisible<2->{\(59\)}&\hspace{0pt}\invisible<2->{\(39\)}&\hspace{0pt}\invisible<2->{\(220\)}&\hspace{0pt}\invisible<2->{\(29\)}\\
\hline
\multicolumn{1}{|r|}{\hspace{0pt}\invisible<2->{Third Dominating rule}}&\hspace{0pt}\invisible<2->{\(56\)}&\hspace{0pt}\hspace{0pt}\invisible<2->{\(39\)}&\hspace{0pt}\invisible<2->{\(220\)}&\hspace{0pt}\hspace{0pt}\invisible<2->{\(29\)}\\
\hline
\multicolumn{1}{|r|}{Experimental class size}&\(56\)&\(39\)&\(213\)&\(29\)\\
\hline
\end{tabularx}
\end{center}
    \caption{Coincidence class number reduction by application of Dominating rules}
    \label{tab:domclasses}
\end{table}
\end{block}
\end{frame}

\subsection{Dominating Pattern Rules}
\label{sub:Dominating pattern Rules}

\begin{frame}{First Dominating rule}
  \begin{block}{Proposition: First Dominating rule}
      \label{prop:dom1}
      \emph{
      Given two mesh patterns \(m_1 =(\sigma, R_1)\) and \(m_2 = (\sigma, R_2)\),
      and a dominating classical pattern \(\pi = (\pi,\emptyset)\) such that
      \(\setsize{\pi} \le \setsize{\sigma} + 1\), the sets \(\av{\{\pi,m_1\}}\) and
      \(\av{\{\pi,m_2\}}\) are coincident if

      \begin{enumerate}
          \item \(R_1 \triangle R_2 = \{(a,b)\}\)
          \item \(\pi \preceq \addp{\sigma}{(a,b)}\)\label{prop:dom1:cont}
      \end{enumerate}}
  \end{block}
\end{frame}

\begin{frame}
  \begin{example}
    The following two patterns are coincident in \(\av{\perm{3,2,1}}\)
    \begin{equation*}
      \begin{tikzpicture}[scale=\picscale]
        \modpattern{}{2,1}{1/0,1/1,1/2,2/2}
    \end{tikzpicture}
        \raisebox{2ex}{\qquad}
    \begin{tikzpicture}[scale=\picscale]
        % \filldraw (2.5,0.5) circle (3pt);
        \modpattern{}{2,1}{1/0,1/1,1/2,2/2,2/0}
    \end{tikzpicture}
    \end{equation*}
  \end{example}
  \begin{corollary}
    All coincidences of classes the form \(\av{\{\perm{3,2,1},\mperm{2,1}{R}}\}\)
    are fully explained by the First Dominating rule.
  \end{corollary}
  \begin{block}{}
    There are \(29\) coincidences of mesh patterns of the form \(\av{\{\perm{3,2,1},\mperm{2,1}{R}}\}\)
  \end{block}
\end{frame}

\begin{frame}{Experimental Results}
  \begin{block}{}
  \begin{table}[htb]
\begin{center}
\begin{tabularx}{\textwidth}{c|Y|Y|Y|Y|}
\cline{2-5}
& \multicolumn{4}{|c|}{Dominating Pattern}\\
\cline{2-5}
& \multicolumn{2}{|c|}{\(\perm{2,3,1}\)}&\multicolumn{2}{|c|}{\(\perm{3,2,1}\)}\\
\cline{2-5}
& \(\perm{1,2}\)&\(\perm{2,1}\)&\(\perm{1,2}\)&\(\perm{2,1}\)\\
\hline
\multicolumn{1}{|r|}{No Dominating rule}&\(220\)&\(220\)&\(220\)&\(220\)\\
\hline
\multicolumn{1}{|r|}{First Dominating rule}&\(85\)&\(43\)&\(220\)&\(29\)\\
\hline
\multicolumn{1}{|r|}{\hspace{0pt}\invisible{Second Dominating rule}}&\hspace{0pt}\invisible{\(59\)}&\hspace{0pt}\invisible{\(39\)}&\hspace{0pt}\invisible{\(220\)}&\hspace{0pt}\invisible{\(29\)}\\
\hline
\multicolumn{1}{|r|}{\hspace{0pt}\invisible{Third Dominating rule}}&\hspace{0pt}\invisible{\(56\)}&\hspace{0pt}\hspace{0pt}\invisible{\(39\)}&\hspace{0pt}\invisible{\(220\)}&\hspace{0pt}\hspace{0pt}\invisible{\(29\)}\\
\hline
\multicolumn{1}{|r|}{Experimental class size}&\(56\)&\(39\)&\(213\)&\(29\)\\
\hline
\end{tabularx}
\end{center}
    \caption{Coincidence class number reduction by application of Dominating rules}
\end{table}
\end{block}
\end{frame}

\begin{frame}{Second Dominating rule}
  \begin{block}{}
    The patterns
    \begin{equation*}
    m_1 = \pattern{}{2,1}{0/0,0/1,0/2,1/0,2/1} \text{ and } m_2 = \pattern{}{2,1}{0/0,0/1,0/2,1/0,1/1,2/1}
    \end{equation*}
    are coincident in \(\av{\perm{2,3,1}}\).
  \end{block}
\end{frame}

\begin{frame}
  \begin{lemma}
    Given a mesh pattern \(m =(\sigma, R)\), where the box \((a,b)\) is not
in \(R\), and a dominating classical pattern \(\pi = (\pi,\emptyset)\) if
\(\pi \preceq \adda{\sigma}{(a,b)}\)\\
,then in any occurrence of \(m\) in a permutation \(\varrho\), the region
corresponding to the box \((a,b)\) can only contain an decreasing
subsequence of \(\varrho\).
  \end{lemma}
\end{frame}

\begin{frame}
  \begin{example}
  Considering \(m_1\) again
  \begin{equation*}
    \begin{tikzpicture}[scale=\picscale]
        {\only<3-4> \fill[pattern=north east lines, pattern color=black!75] (1,1) rectangle +(1,1);}
        \modpattern{}{2,1}{0/0,0/1,0/2,1/0,2/1}
        {\only<2> \filldraw (1.66,1.33) circle (3pt);}
        {\only<2> \filldraw (1.33,1.66) circle (3pt);}
        {\only<3> \filldraw (2.66,0.33) circle (3pt);}
        {\only<3> \filldraw (2.33,0.66) circle (3pt);}
    \end{tikzpicture}
  \end{equation*}
  {\only<1-3> \vphantom{This is \(m_2\).}}
  {\only<4-> {}}
  \end{example}
\end{frame}

\begin{frame}
  \begin{block}{Proposition: Second Dominating rule}
    \label{prop:dom2}
    \emph{
    Given two mesh patterns \(m_1 =(\sigma, R_1)\) and \(m_2 = (\sigma, R_2)\),
    and a dominating classical pattern \(\pi = (\pi,\emptyset)\) such that
    \(\setsize{\pi} \le \setsize{\sigma} + 2\), the sets \(\av{\{\pi,m_1\}}\) and
    \(\av{\{\pi,m_2\}}\) are coincident if

    \begin{enumerate}
        \item \(R_1 \triangle R_2 = \{(a,b)\}\)
        \item   \begin{enumerate}
                \item\label{prop:dom2:condc} \(\pi \preceq \adda{\sigma}{(a,b)}\) and
                        \begin{enumerate}
                            \item \((a+1,b) \in \sigma\) and \((a+1,b-1)\notin R\) and \\
                                \((x,b-1)\in R \implies (x,b) \in R \) (where \(x\neq a,a+1\)) and\\
                                  \((a+1,y)\in R \implies (a,y) \in R\) (where \(y\neq b-1,b\)).
                            \item \(\dots\)
                        \end{enumerate}
                    \item \(\dots\)\begin{enumerate}
                            \item \(\dots\)
                            \item \(\dots\)
                          \end{enumerate}
                \end{enumerate}
    \end{enumerate}
    }
  \end{block}
\end{frame}

\begin{frame}
  \begin{example}
    \begin{figure}[ht]
\centering
\begin{tikzpicture}[scale = 0.75]

     % Make the grid
     \draw[step=1cm,gray!50] (2/3,2/3) grid (5+1/3,5+1/3);

     % Make black cross in the middle
     \draw[-,very thick] (3,2/3) -- (3,5+1/3);
     \draw[-,very thick] (2/3,3) -- (5+1/3,3);
     \draw[-,very thick,gray] (4,2) -- (2,2) -- (2,4) -- (4,4) -- (4,2);

     \fill[black] (3,3) circle (0.1);
     \fill[red] (2,3) circle (0.1);

     % Draw the dots on the side
     \foreach \x in {1,2,4,5}{
         \draw[dotted, gray!50] (\x, 0) -- (\x, 2/3);
         \draw[dotted, gray!50] (\x, 5+1/3) -- (\x, 6);
         \draw[dotted, gray!50] (0,\x) -- (2/3,\x);
         \draw[dotted, gray!50] (5+1/3,\x) -- (6,\x);
        }
     % Draw black dots next to the cross
     \draw[dotted, black, very thick] (3, 0) -- (3, 2/3);
     \draw[dotted, black, very thick] (3, 5+1/3) -- (3, 6);
     \draw[dotted, black, very thick] (0, 3) -- (2/3, 3);
     \draw[dotted, black, very thick] (5+1/3, 3) -- (6, 3);

     \draw[->-=.5,thick] (3,3) to[bend right=15](2,4);

     \draw[->, thick] (5,3) -- (5,4);
     \draw[->, thick] (1,3) -- (1,4);
     \draw[->, thick] (3,5) -- (2,5);
     \draw[->, thick] (3,1) -- (2,1);

     \node [black,right] at (6,3) {$b$};
     \node [black,below] at (2,0) {$a$};

\end{tikzpicture}
    \caption{If the conditions of The Second Dominating rule are satisfied the box
    \((a,b)\) can be shaded.}
\end{figure}
  \end{example}
\end{frame}

\begin{frame}
  \begin{corollary}
    All coincidences of classes the form \(\av{\{\perm{2,3,1},\mperm{2,1}{R}}\}\)
    are fully explained by applying the First Dominating rule, then applying the
    Second Dominating rule.
  \end{corollary}
  \begin{block}{}
    There are \(39\) coincidences of mesh patterns of the form
    \(\av{\{\perm{2,3,1},\mperm{2,1}{R}}\}\)
  \end{block}
\end{frame}

\begin{frame}{Experimental Results}
  \begin{block}{}
  \begin{table}[htb]
\begin{center}
\begin{tabularx}{\textwidth}{c|Y|Y|Y|Y|}
\cline{2-5}
& \multicolumn{4}{|c|}{Dominating Pattern}\\
\cline{2-5}
& \multicolumn{2}{|c|}{\(\perm{2,3,1}\)}&\multicolumn{2}{|c|}{\(\perm{3,2,1}\)}\\
\cline{2-5}
& \(\perm{1,2}\)&\(\perm{2,1}\)&\(\perm{1,2}\)&\(\perm{2,1}\)\\
\hline
\multicolumn{1}{|r|}{No Dominating rule}&\(220\)&\(220\)&\(220\)&\(220\)\\
\hline
\multicolumn{1}{|r|}{First Dominating rule}&\(85\)&\(43\)&\(220\)&\(29\)\\
\hline
\multicolumn{1}{|r|}{Second Dominating rule}&\(59\)&\(39\)&\(220\)&\(29\)\\
\hline
\multicolumn{1}{|r|}{\hspace{0pt}\invisible{Third Dominating rule}}&\hspace{0pt}\invisible{\(56\)}&\hspace{0pt}\invisible{\(39\)}&\hspace{0pt}\invisible{\(220\)}&\hspace{0pt}\invisible{\(29\)}\\
\hline
\multicolumn{1}{|r|}{Experimental class size}&\(56\)&\(39\)&\(213\)&\(29\)\\
\hline
\end{tabularx}
\end{center}
    \caption{Coincidence class number reduction by application of Dominating rules}
\end{table}
\end{block}
\end{frame}

\begin{frame}{Third Dominating rule}
  \begin{block}{}
    The patterns
    \begin{equation*}
    m_1 = \pattern{}{1,2}{0/2,2/0,2/1} \text{ and } m_2 = \pattern{}{1,2}{0/2,1/0,2/0,2/1}
    \end{equation*}
    are coincident in \(\av{\perm{2,3,1}}\).
    Neither of the previous two rules explain this.
  \end{block}
\end{frame}

\begin{frame}
  \begin{example}
    \begin{equation*}
    \only<1>{\vphantom{\pattern{2,3}{2,1,3}{0/3,1/3,2/0,2/1,3/0,3/1,3/2}}\pattern{}{1,2}{0/2,2/0,2/1}}
    \only<2>{\pattern{}{2,1,3}{0/3,2/0,2/1,3/0,3/1,3/2}}
    \only<3>{\pattern{}{2,1,3}{0/3,1/3,2/0,2/1,3/0,3/1,3/2}}
    \only<4>{\pattern{2,3}{2,1,3}{0/3,1/3,2/0,2/1,3/0,3/1,3/2}}
    \only<5->{\vphantom{\pattern{2,3}{2,1,3}{0/3,1/3,2/0,2/1,3/0,3/1,3/2}}\pattern{}{1,2}{0/2,1/0,2/0,2/1}}
  \end{equation*}
  \end{example}
\end{frame}

\begin{frame}{}
  \begin{block}{Proposition: Third Dominating rule}
      \emph{
      Given two mesh patterns \(m_1 =(\sigma, R_1)\) and \(m_2 = (\sigma, R_2)\),
      and a dominating classical pattern \(\pi = (\pi,\emptyset)\), the sets
      \(\av{\{\pi,m_1\}}\) and \(\av{\{\pi,m_2\}}\) are coincident if
      \begin{enumerate}
          \item \(R_1 \triangle R_2 = \{(a,b)\}\)
          \item\label{prop:dom3:condocc} \(\addp[D]{\mperm{\sigma}{R_1}}{(a,b)}\)
              where \(D\in\{N,E,S,W\}\)
              is coincident with a mesh pattern containing an occurrence of
              \(\mperm{\sigma}{R_2}\) as a subpattern.
      \end{enumerate}
      }
  \end{block}
\end{frame}

\begin{frame}
  \begin{corollary}
    All coincidences of classes the form \(\av{\{\perm{2,3,1},\mperm{1,2}{R}}\}\)
    are fully explained by applying the First Dominating rule, the
    Second Dominating rule, and then the Third Dominating rule.
  \end{corollary}
  \begin{block}{}
    There are \(56\) coincidences of mesh patterns of the form
    \(\av{\{\perm{2,3,1},\mperm{2,1}{R}}\}\)
  \end{block}
\end{frame}

\begin{frame}{Experimental Results}
  \begin{block}{}
  \begin{table}[htb]
\begin{center}
\begin{tabularx}{\textwidth}{c|Y|Y|Y|Y|}
\cline{2-5}
& \multicolumn{4}{|c|}{Dominating Pattern}\\
\cline{2-5}
& \multicolumn{2}{|c|}{\(\perm{2,3,1}\)}&\multicolumn{2}{|c|}{\(\perm{3,2,1}\)}\\
\cline{2-5}
& \(\perm{1,2}\)&\(\perm{2,1}\)&\(\perm{1,2}\)&\(\perm{2,1}\)\\
\hline
\multicolumn{1}{|r|}{No Dominating rule}&\(220\)&\(220\)&\(220\)&\(220\)\\
\hline
\multicolumn{1}{|r|}{First Dominating rule}&\(85\)&\(43\)&\(220\)&\(29\)\\
\hline
\multicolumn{1}{|r|}{Second Dominating rule}&\(59\)&\(39\)&\(220\)&\(29\)\\
\hline
\multicolumn{1}{|r|}{Third Dominating rule}&\(56\)&\(39\)&\(220\)&\(29\)\\
\hline
\multicolumn{1}{|r|}{Experimental class size}&\(56\)&\(39\)&\(213\)&\(29\)\\
\hline
\end{tabularx}
\end{center}
    \caption{Coincidence class number reduction by application of Dominating rules}
\end{table}
\end{block}
\end{frame}

\subsection{Special Cases}
\label{sub:Special Cases}
\begin{frame}{Special Cases}
  There are \(7\) coincidences of the form \(\av{\perm{3,2,1},m}\) that are not
  explained by the rules.
  \begin{example}
    \begin{equation*}
      m_1 = \pattern{}{1,2}{0/1,0/2,2/0,2/1} \text{ and } m_2 = \pattern{}{1,2}{0/1,0/2,1/1,2/0,2/1}
    \end{equation*}
    This coincidence is explained by mathematical induction on the number of
    points in the region corresponding to the middle box.
  We call this number \(n\).
  \end{example}
\end{frame}
\begin{frame}
  \begin{example}
    \begin{description}
    \only<1-2>{\item [Base Case \((n=0)\):] The base case holds since we can freely shade the
        box if it contains no points.}
    \only<2>{\item [Inductive Hypothesis \((n=k)\):] Suppose that we can find an occurrence
        of the second pattern if we have an occurrence of the first with \(k\) points
        in the middle box.}
    \only<3>{\item [Inductive Step \((n=k+1)\)] Suppose that we have \((k+1)\) points in
        the middle box. Choose the bottom most point in the middle box, this gives the
        mesh pattern
     \begin{equation*}
        \begin{tikzpicture}[scale=0.5]
            \modpattern[5]{}{1,2,3}{0/1,0/2,0/3,1/1,2/1,3/0,3/1,3/2}
            \draw (2.5,0.5) node {\(X\)};
        \end{tikzpicture}
    \end{equation*}
        Now we need to consider the box labelled \(X\). If this box is empty then
        we have an occurrence of \(m_2\) and are done.}
        \only<4>{\item [Inductive Step \((n=k+1)\) (cont.)] If this box contains
        any points then we gain some extra shading on the mesh pattern
        due to the dominating pattern
     \begin{equation*}
        \begin{tikzpicture}[scale=0.5]
            \modpattern[5]{2,3}{1,2,3}{0/1,0/2,0/3,1/1,1/2,1/3,2/1,3/0,3/1,3/2}
            \filldraw (2.5,0.5) circle (4 pt);
        \end{tikzpicture}
    \end{equation*}
        The two highlighted points form an occurrence of \(m_1\) with \(k\)
        points in the middle box, and thus by the Inductive Hypothesis we
        are done.}
\end{description}
  \end{example}
\end{frame}
